% AIrxiv Primer
% Generated 2026-03-03
\documentclass[11pt]{article}

\usepackage[a4paper,margin=1in]{geometry}
\usepackage{amsmath,amssymb,amsthm}
\usepackage{mathtools}
\usepackage{hyperref}
\usepackage{enumitem}
\usepackage{microtype}

\title{From Markov Exchangeability to Mixtures of Markov Chains:\\
A Conceptual Primer for a Lean Formalization (Fortini--Diaconis--Freedman Route)}
\author{AIrxiv note (conceptual; implementation-oriented)}
\date{2026-03-03}

\newcommand{\N}{\mathbb{N}}
\newcommand{\Fin}{\mathrm{Fin}}
\newcommand{\Law}{\mathcal{L}}
\newcommand{\cyl}{\mathrm{cyl}}
\newcommand{\E}{\mathbb{E}}
\newcommand{\Prob}{\mathbb{P}}

\begin{document}
\maketitle

\begin{abstract}
This note is a didactic, strictly conceptual primer for a formalization of the Markov de Finetti / Fortini representation theorem in a proof assistant (Lean).
The key idea is to state the \emph{public} theorem with classical hypotheses (Markov exchangeability plus recurrence) and to treat low-level ``constructor payloads''
(e.g.\ measurability side-conditions, start-restricted factorization, and stepwise induction gadgets) as \emph{internal} lemmas rather than as user-facing assumptions.

The central lesson is that the remaining gap in many formal developments is not ``more measure theory'': it is a \emph{joint coherence principle}
about the whole successor-state matrix, which cannot be reconstructed from per-row statements alone.
We summarize the canonical Fortini route (successor matrix partial exchangeability \(\Rightarrow\) mixture of Markov chains) and relate it to a second,
more constructive route based on certified finite combinatorics (``BEST''-style quantitative bounds), clarifying where verified algorithms help and where they do not.
\end{abstract}

\tableofcontents

\section{What problem is being formalized?}

Fix a finite state space \(S = \Fin(k)\). A stochastic process \(X_0,X_1,\dots\) with values in \(S\) is:

\begin{itemize}[leftmargin=2em]
\item \emph{a (time-homogeneous) Markov chain} if, conditional on \(X_n\), the future depends on the past only via \(X_n\), with a transition matrix \(P\).
\item \emph{a mixture of Markov chains} if there exists a probability measure \(\pi\) on transition parameters \(\theta\) such that the law of \(X\) is obtained by first sampling \(\theta \sim \pi\), then running the Markov chain with parameter \(\theta\).
\end{itemize}

A Markov de Finetti theorem characterizes which laws of \(X\) admit such a representation by an \emph{invariance property} on observable finite prefixes.

\paragraph{Classical statement (informal).}
Under a suitable notion of \emph{Markov exchangeability} and a recurrence condition, the law of \(X\) is a mixture of Markov chains.

A major reference for the successor-matrix formulation is Fortini--Ladelli--Petris--Regazzini (2002), which proves an equivalence:
a law is a mixture of Markov chains iff the \emph{successor-state matrix} (defined below) is invariant under finite permutations within rows,
and under recurrence this row-wise partial exchangeability is equivalent to Markov exchangeability
in the sense of Diaconis--Freedman. \cite{FortiniEtAl2002}

\section{Two invariance principles: Markov exchangeability vs successor-matrix PE}

\subsection{Markov exchangeability}

Markov exchangeability is an invariance of the \emph{prefix law} under permutations of time indices that preserve the relevant Markov ``evidence'' of a finite word.
In informal combinatorial terms: two finite paths with the same start state and the same transition counts should be assigned the same probability.

Different formalizations package this differently (e.g.\ via equivalence relations on words or via ``evidence-preserving'' bijections between cylinder events).
What matters conceptually is that Markov exchangeability is a \emph{symmetry of the observed path}.

\subsection{Successor-state matrix and partial exchangeability}

Define, for each \(i \in S\) and \(n \in \N\), the \emph{successor after the \(n\)-th visit to \(i\)}:
\[
V_{i,n} := \text{the value of \(X\) immediately after the \(n\)-th time \(X\) visits state \(i\)}.
\]
The array \(V = (V_{i,n})_{i \in S,\, n \in \N}\) is the \emph{successor-state matrix}.

\paragraph{Row-wise partial exchangeability (PE).}
The law of \(V\) is invariant under applying an arbitrary finite permutation \(\sigma_i\) to the indices \(n\) within each row \(i\), jointly across finitely many coordinates:
for any finite set of queried entries, permuting visit numbers independently within each row does not change the joint distribution.

This is \emph{not} a purely per-row condition. It is a statement about the \emph{joint law} of the entire matrix \(V\), allowing correlations between rows.
Fortini et al.\ show: mixture-of-Markov-laws \(\Leftrightarrow\) row-wise PE of \(V\) (plus mild technicalities for general state spaces). \cite{FortiniEtAl2002}

\subsection{Why recurrence appears (and why ``strong recurrence'' may show up in Lean)}

A technical subtlety: for some sample paths, a state \(i\) may be visited only finitely many times, so \(V_{i,n}\) eventually becomes undefined.
Fortini et al.\ handle this by adjoining a dummy value (often denoted \(\partial\)) to each row so the matrix is always total. \cite{FortiniEtAl2002}

In a formalization that keeps the state space fixed as \(\Fin(k)\), one often replaces this with a hypothesis ensuring that the \(n\)-th visit time exists whenever it is needed.
This yields a \emph{strong recurrence} (or \emph{infinite visits}) assumption:
\[
\text{for each \(i\), on almost every path that ever visits \(i\), all \(n\)-th visit times exist.}
\]
This is logically close to recurrence but can be more convenient than introducing an explicit dummy state into the type.

\section{The real crux in formalization: joint coherence vs per-row data}

\subsection{A useful internal factorization: row kernels}

A common formal tactic is to construct, for each anchor state \(i\), a \emph{row kernel}
\[
K_i(r) \in \mathrm{Prob}(S),
\]
where \(r\) is (some encoding of) the row-process history and \(K_i(r)\) predicts successors from state \(i\).
Row-wise PE suggests that, conditional on a latent directing object, successors within row \(i\) are i.i.d.\ with law \(K_i(r)\).

This row-kernel viewpoint is powerful for building the final integral representation, but it also introduces a trap:

\begin{quote}
Per-row factorization statements (even if true for every row) do \emph{not} determine the joint law across rows.
The missing piece is a \emph{cross-row coherence} principle ensuring that the row kernels come from a \emph{single latent transition mechanism}.
\end{quote}

\subsection{Cross-anchor identities as a coherence surrogate}

A pragmatic way to express cross-row coherence in code is via a \emph{cross-anchor product identity}:
an identity equating the probability of a cylinder event with an integral of a product of row-kernel evaluations,
uniformly across choices of the ``anchor'' used to parameterize the construction.

Conceptually, this identity is a shadow of the statement:
\[
\exists \text{ latent transition matrix } \Theta \text{ such that } V_{i,n} \text{ are conditionally i.i.d.\ within each row and coherent across rows.}
\]

In many developments, a stepwise recursion hypothesis (call it ``\(\mathsf{hstep}\)'') appears, used to lift a length-two base case to arbitrary lengths.
This is mathematically legitimate as an \emph{internal lemma}, but it should not be a \emph{public} assumption:
in the literature, the correct public assumption is successor-matrix PE (or Markov exchangeability + recurrence), from which such internal recursion can be derived.

\subsection{Negative lesson: ``bad row kernels''}

It is straightforward to build examples where row-wise-looking ingredients exist but cross-anchor coherence fails:
choose row kernels that do not arise from any single transition matrix, and cross-anchor identities can break.
This demonstrates that a formal proof must import some \emph{joint} hypothesis---successor-matrix PE is the canonical one.

\section{A clean theorem hierarchy (recommended)}

The following hierarchy keeps user-facing statements classical and pushes proof plumbing into internal lemmas.

\subsection{Public theorem (literature-facing)}

\begin{quote}
\textbf{Markov de Finetti for Markov chains (Fortini--Diaconis--Freedman route).}
If a discrete-valued process is Markov exchangeable and recurrent (or strongly recurrent, depending on encoding),
then its law is a mixture of Markov chains.
\end{quote}

The best practice is to present this as:
\[
\text{(Markov exchangeability + recurrence)} \Rightarrow \text{(successor matrix PE)} \Rightarrow \text{(mixture of Markov chains).}
\]
Diaconis--Freedman (1980) is a classical source for the Markov-chain de Finetti theorem, and Fortini et al.\ (2002) cleanly isolates the successor-matrix PE condition. \cite{DiaconisFreedman1980,FortiniEtAl2002}

\subsection{Canonical intermediate definitions}

\begin{itemize}[leftmargin=2em]
\item \textbf{Strong recurrence} (if no dummy state is used): ensures the successor matrix is total almost surely on relevant events.
\item \textbf{SuccessorMatrixPE}: joint invariance under finite within-row permutations.
\item \textbf{RowSuccessorMatrixInvariance}: a derived, often more convenient invariance principle tailored to the row-successor process used in the implementation.
\end{itemize}

\subsection{Internal lemmas (implementation helpers)}

These should be marked clearly internal:
\begin{itemize}[leftmargin=2em]
\item Measurability obligations for \(\Fin(1)\) ``pi'' measures (often derivable from evaluation measurability on atoms).
\item Start-restricted row-law factorization (a conditioning/disintegration byproduct).
\item Stepwise induction gadgets (\(\mathsf{hstep}\)-style lemmas): useful for proofs, not for theorem statements.
\end{itemize}

\section{Where verified algorithms help (and where they do not)}

A second proof approach uses \emph{quantitative finite} control (WR/WOR transport bounds, certificates, checkers).
This is valuable, but it targets a different kind of claim.

\subsection{Strengths}

Verified algorithms shine when the obligation is:
\begin{itemize}[leftmargin=2em]
\item finite combinatorics (counting / carrier equivalence / finite-dimensional reductions),
\item explicit inequalities with rational certificates,
\item automatic falsification of incorrect finite claims (counterexample checkers).
\end{itemize}

This makes them perfect for:
\begin{itemize}[leftmargin=2em]
\item pruning false ``shortcut'' lemmas early,
\item validating finite pattern-rate bounds used in a constructive route,
\item supporting a \emph{quantitative} variant of the representation theorem (stronger hypotheses, more computable conclusions).
\end{itemize}

\subsection{Limitations}

They do \emph{not} replace the analytic core:
\begin{itemize}[leftmargin=2em]
\item almost-everywhere statements on infinite paths,
\item disintegration/conditioning arguments,
\item construction of a latent directing measure from a joint invariance principle.
\end{itemize}

Thus, the most satisfying integration is:
\begin{quote}
Use checkers/certificates to \emph{guard} finite combinatorial lemmas and to build constructive quantitative variants,
while proving the canonical Fortini bridge (Markov exchangeability + recurrence \(\Rightarrow\) successor-matrix PE) by a human-readable combinatorial argument.
\end{quote}

\section{How to keep the codebase (and future AI sessions) clean}

A highly effective organizational pattern is a three-tier split:
\begin{description}[leftmargin=2.8em,style=nextline]
\item[Canonical] theorems and definitions intended for users: the public theorem, successor-matrix PE, recurrence assumptions, and the minimal composition chain.
\item[Internal] builder payloads, measurability plumbing, start-restriction utilities, intermediate constructors; explicitly documented as internal scaffolding.
\item[Legacy/Archive] historical proof attempts, deprecated interfaces, and exploratory lemmas; physically separated so they do not distract new development sessions.
\end{description}

In practice:
\begin{itemize}[leftmargin=2em]
\item ensure the main exported file imports only \texttt{Core + Canonical};
\item move deprecated or exploratory interfaces into a \texttt{Legacy} module;
\item move historical experiments into an \texttt{Archive} tree with a short README stating ``not imported by active proofs.''
\end{itemize}

\section{A suggested ``completion roadmap'' for the canonical route}

Let \(\mu\) denote a prefix law and \(P\) a path measure extending it.

\begin{enumerate}[leftmargin=2em]
\item \textbf{(Main gap)} Prove:
\[
\text{Markov exchangeability + (strong) recurrence} \Rightarrow \text{SuccessorMatrixPE}.
\]
This is primarily a finite-combinatorics argument using prefix-carrier equivalences and cylinder decompositions.
\item From SuccessorMatrixPE, extract:
\[
\exists \text{ row kernels with evaluation measurability } + \text{ row-successor invariance}.
\]
\item Combine with the already-established ``downstream'' implication:
\[
\text{row-successor invariance} \Rightarrow \text{cross-anchor identity} \Rightarrow \text{Markov-mixture representation}.
\]
\end{enumerate}

The key conceptual discipline is to keep \(\mathsf{hstep}\)-style assumptions out of the public story:
if a stepwise induction is needed internally, it should be derived from successor-matrix PE, not postulated.

\section{Conclusion}

A mathematically satisfying formalization of the Markov de Finetti theorem for Markov chains is best obtained by:
\begin{itemize}[leftmargin=2em]
\item stating the \emph{public} theorem using classical hypotheses (Markov exchangeability + recurrence/strong recurrence);
\item making successor-matrix PE the primary intermediate invariant;
\item treating row-kernel constructors and stepwise recursion as internal proof structure;
\item using verified algorithms to certify finite obligations and to prevent false combinatorial shortcuts, while keeping the analytic bridge human-readable.
\end{itemize}

\bigskip
\noindent\textbf{Acknowledgment.} This note is intentionally conceptual and implementation-oriented; it does not attempt to reproduce full proofs from the cited literature, but rather to clarify how their hypothesis structure maps cleanly into a modern mechanized proof development.

\bibliographystyle{plain}
\begin{thebibliography}{9}

\bibitem{DiaconisFreedman1980}
P.~Diaconis and D.~Freedman.
\newblock De {F}inetti's theorem for {M}arkov chains.
\newblock \emph{The Annals of Probability}, 8(1):115--130, 1980.
\newblock Available via Project Euclid: \url{http://projecteuclid.org/euclid.aop/1176994828}.

\bibitem{FortiniEtAl2002}
S.~Fortini, L.~Ladelli, G.~Petris, and E.~Regazzini.
\newblock On mixtures of distributions of {M}arkov chains.
\newblock \emph{Stochastic Processes and their Applications}, 100(1--2):147--165, 2002.
\newblock DOI: \url{https://doi.org/10.1016/S0304-4149(02)00093-5}.

\bibitem{FortiniPetrone2017}
S.~Fortini and S.~Petrone.
\newblock Predictive characterization of mixtures of {M}arkov chains.
\newblock \emph{Bernoulli}, 2017.
\newblock DOI: \url{http://dx.doi.org/10.3150/15-BEJ787}.

\end{thebibliography}

\end{document}
